\section{Chapter Workshop}

\begin{workedexample}{Bayesian classification}
A test has sensitivity $0.95$ and specificity $0.90$. Suppose the prevalence of the condition is $0.01$. If $D$ is the condition and $+$ is a positive result, then
\[
\mathbb{P}(D\mid +)
=\frac{\mathbb{P}(+\mid D)\mathbb{P}(D)}
{\mathbb{P}(+\mid D)\mathbb{P}(D)+\mathbb{P}(+\mid D^c)\mathbb{P}(D^c)}.
\]
Substitution gives
\[
\mathbb{P}(D\mid +)=\frac{0.95\cdot0.01}{0.95\cdot0.01+0.10\cdot0.99}\approx0.0876.
\]
A highly sensitive test can still have a modest positive predictive value when the base rate is small.
\end{workedexample}

\newpage

\begin{applicationbox}{Monte Carlo integration}
If $U_1,\dots,U_n$ are independent uniform samples on $[0,1]$, then
\[
    \widehat I_n=\frac1n\sum_{k=1}^n f(U_k)
\]
estimates $\int_0^1f(x)\,dx$. The law of large numbers gives consistency, while the central limit theorem gives an approximate error scale
\[
    \frac{\sqrt{\operatorname{Var}(f(U))}}{\sqrt n}.
\]
The method extends to high-dimensional integrals where deterministic grids become prohibitively expensive.
\end{applicationbox}

\begin{applicationbox}{Markov models and ranking}
A finite Markov chain evolves a probability vector by multiplication with a stochastic matrix. Stationary distributions describe long-run occupancy and underlie models of queues, population movement, and ranking on directed networks. Linear algebra determines invariant distributions; probability determines how sample paths approach them.
\end{applicationbox}

\begin{chapterexercises}
    \item Two fair dice are rolled. Compute the conditional probability that the sum is $8$ given that at least one die is $3$.
    \item Let $X\sim\mathrm{Binomial}(n,p)$. Derive $\mathbb{E}X$ and $\operatorname{Var}(X)$ using indicator variables.
    \item If $X$ has density $f(x)=2x$ on $[0,1]$, find its distribution function, mean, and variance.
    \item Give an example showing that pairwise independence need not imply mutual independence.
    \item Prove Markov's inequality and derive Chebyshev's inequality from it.
    \item Explain the difference between convergence in probability and almost-sure convergence.
    \item A Markov chain has transition matrix $P=\begin{pmatrix}0.8&0.2\\0.3&0.7\end{pmatrix}$. Find its stationary distribution.
    \item For a Poisson process of rate $\lambda$, derive the distribution of the waiting time to the first arrival.
\end{chapterexercises}

\begin{selectedhints}
    \item The conditioning event has $11$ ordered outcomes; two of them, $(3,5)$ and $(5,3)$, have sum $8$, so the probability is $2/11$.
    \item Write $X=I_1+\cdots+I_n$ with independent Bernoulli indicators.
    \item $F(x)=x^2$ on $[0,1]$; $\mathbb{E}X=2/3$ and $\operatorname{Var}(X)=1/18$.
    \item Let two fair bits $A,B$ be independent and set $C=A\oplus B$.
    \item Apply Markov to the nonnegative random variable $(X-\mu)^2$.
    \item Almost-sure convergence controls entire sample paths outside one null set; convergence in probability controls the probability of a deviation separately for each $n$.
    \item Solve $\pi P=\pi$ with $\pi_1+\pi_2=1$; the answer is $(3/5,2/5)$.
\end{selectedhints}
